\documentclass[reqno]{exam}

\usepackage{amssymb, amsmath, amsthm, stmaryrd}
\usepackage{fullpage, parskip}
\usepackage{graphicx}
\usepackage{enumitem}

% Improved spacing
\usepackage[top=1in, bottom=0.8in, left=1in, right=1in]{geometry}


\title{\normalsize{
  JHED: \underline{\hspace{4cm}} \hspace{1em} Name: \underline{\hspace{9cm}}}
   }
\author{}
\date{}

% Custom spacing for better page utilization
\renewcommand{\baselinestretch}{0.95}

\begin{document}
\maketitle
\vspace{-2cm} % Reduce space after title

\begin{questions}

\question 
Find the line of best fit for the following data points: $(-1, 0), (0, 4), (1, 2), (2, 4)$.
\begin{solution}
Let
\[
A=\begin{bmatrix}
-1 & 1\\
0 & 1\\
1 & 1\\
2 & 1
\end{bmatrix},
\quad
\mathbf y=\begin{bmatrix}
0\\4\\2\\4
\end{bmatrix},
\quad
\text{and } \mathbf y \approx A\begin{bmatrix}m\\c\end{bmatrix}.
\]
The normal equations are
\[
A^TA\begin{bmatrix}m\\c\end{bmatrix}=A^T\mathbf y.
\]
Compute:
\[
A^TA=\begin{bmatrix}6&2\\2&4\end{bmatrix},
\qquad
A^T\mathbf y=\begin{bmatrix}2\\6\end{bmatrix}.
\]
So
\[
\begin{bmatrix}m\\c\end{bmatrix}
=(A^TA)^{-1}A^T\mathbf y
=\begin{bmatrix}-\frac15\\[2pt]\frac85\end{bmatrix}.
\]
Hence, the line of best fit is
\[
y=-\frac15x+\frac85.
\]
\end{solution}

\newpage
\question
\begin{parts}
  \part Find two different orthogonal matrices in $Q_1$ and $Q_2$ that send the vector $\begin{bmatrix}
    1 \\ 2
  \end{bmatrix}$ to $\begin{bmatrix}
    -2 \\ 1
  \end{bmatrix}$.
  
  \begin{solution}
    
  \end{solution}

  \part Find all vectors $\vec{v}$ such that $Q_1 \vec{v} = \vec{v}$.
  
  \begin{solution}
  $Q_1$ is a rotation by $\pi/4$, which has no nonzero eigenvector with eigenvalue 1.

  $\vec{v} = \begin{bmatrix} 0 \\ 0 \end{bmatrix}$ only.
  \end{solution}

  \part Find all vectors $\vec{v}$ such that $Q_2 \vec{v} = \vec{v}$.
  
  \begin{solution}
  $Q_2$ is a reflection. Fixed vectors lie along the midpoint of $\vec{a}$ and $\vec{b}$:
  
  $\vec{a} + \vec{b} = \begin{bmatrix} 1 \\ 1 \end{bmatrix} + \begin{bmatrix} 0 \\ \sqrt{2} \end{bmatrix} = \begin{bmatrix} 1 \\ 1 + \sqrt{2} \end{bmatrix}$

  $\vec{v} = t\begin{bmatrix} 1 \\ 1 + \sqrt{2} \end{bmatrix}$ for $t \in \mathbb{R}$.
  \end{solution}
\end{parts}
\end{questions}
\end{document}